\fichatitulo{43 · METODOLOGIA}{MIS Quarterly · 2004}{Design Science em Sistemas de Informação}
{\small Hevner et al. \textit{Design Science in Information Systems Research}. MIS Quarterly · 2004. \href{https://ailab-ua.github.io/courses/MIS611D/hevner-2004-design_science_in_is_research.pdf}{Fonte consultada}.\par}

\campo{Problema e objetivo}
Como produzir conhecimento por meio da construção e avaliação de artefatos?

\campo{Principal contribuição · síntese interpretativa}
Articula relevância do problema e rigor da pesquisa em sete diretrizes para Design Science.

\campo{Método / natureza da evidência}
Ensaio metodológico com quadro conceitual e aplicação ilustrativa a três pesquisas; o recorte focal concentra-se nas diretrizes.

\campo{Resultado ou argumento principal}
Exige artefato viável, problema relevante, avaliação, contribuição verificável, rigor, busca de soluções e comunicação apropriada.

\campo{Excertos-chave · seleção literal}
\begin{itemize}
\excerto{1}{Design-science research must produce a viable artifact}{Tabela 1, diretriz 1.}
\excerto{2}{The utility, quality, and efficacy of a design artifact must be rigorously demonstrated}{Diretriz 3.}
\end{itemize}

\campo{Limitações e análise crítica}
Análise da ficha: construir um protótipo não basta para demonstrar contribuição científica. Critérios de avaliação devem decorrer do problema e do contexto organizacional, sem tratar as diretrizes como formulário mecânico.

\campo{Aproveitamento na dissertação · proposta}
Metodologia: relacionar o artefato RAG, o problema parlamentar, as decisões de projeto e a evidência de utilidade.

\registro{Fonte consultada nesta preparação: Resumo, introdução, tabela 1 e discussão das diretrizes 1–3; exemplares não fichados. Chave: \texttt{hevnerEtAl2004designScience}.}
